\section{Expansion and periodic lines in the support}
\label{exp:section}

We retain the nonzero stationary current $\mu$, its closed support
$\Omega\subset\partial^2P$, and the invariant conull set $\mathcal U$
from the preceding sections. Our purpose here is to exclude periodic
lines from $\Omega$. This is stronger than atomlessness: a line can belong
to the support while having measure zero. Only after proving this stronger
assertion will we extend the trajectory construction from $\mathcal U$
to every line of $\Omega$.

The obstruction is a possible carrier line in the support despite its
zero measure. We will expand a positive-measure cylinder near such a
line along a fixed return. If the whole cylinder follows the return
forever, its translated images concentrate on the zero-mass carrier
set. If a line leaves the prescribed return domain, it follows a
competing window; the cyclic intersection of the two window stabilizers
forces a periodic middle in their long common segment. Lemma~\ref{exp:finite-failure}
records this alternative quantitatively. We first use it to produce
periodic support from carrier support, then use whole-period pullback
to rule out periodic support itself.

The window results provide the finite data for this argument. For each
fixed shape there are finitely many representative cores of noncyclic
window stabilizers; only finitely many actual windows have their hulls
meeting a fixed ball. Distinct actual windows of that shape have at
most cyclic common stabilizer, and an irrational endpoint determines
at most one window (Proposition~\ref{stat:windows} and
Lemma~\ref{stat:window-finiteness}). We will also use the fixed small
subgraph of the selected cores $\Gamma_n$ from
Definition~\ref{car:small-graph}.

Let $M\le P$ represent one of the maximal noncyclic carrier classes.
Proposition~\ref{car:return} and Lemma~\ref{car:expansion} supply
an automorphism $\beta$ of $M$, a number $\lambda>4$, and fixed forward
and backward window transfers $\Theta_+$ and $\Theta_-$ with domains
\[
 \partial^2\Lambda S_+\quad\hbox{and}\quad
 \partial^2\Lambda S_-,\qquad M\le S_+,S_-\le P.
\]
Here and below $\partial^2\Lambda K$ denotes the ordered pairs of
distinct points of the limit set of $K$. On the boundary of $M$ these
transfers are respectively $\beta$ and $\beta^{-1}$. They are induced
by the actual paths to the positive and negative powers of the return
frame. In an $M$-basis,
\begin{equation}
 \lambda |w|_M\le
 \max\{|\beta(w)|_M,|\beta^{-1}(w)|_M\}
 \qquad(w\in M).
 \label{exp:based-inequality}
\end{equation}
The same inequality holds for cyclic lengths. Moreover,
Proposition~\ref{car:zero-mass} gives
\begin{equation}
 \mu(\partial^2\Lambda M)=0.
 \label{exp:zero-carrier-mass}
\end{equation}
All transfers preserve $\mu$ on their domains in $\mathcal U$ and
take those domains into $\mathcal U$. No assertion about their action
on arbitrary support lines is being assumed.

We call a line of the $P$-Cayley tree an \emph{$M$-line} if its two
endpoints belong to $\Lambda M$. It corresponds to a unique line of
the $M$-Cayley tree under the boundary embedding. A periodic line is
the axis of a nontrivial element. Its primitive translation is a
generator of its cyclic stabilizer; this does not mean a member of a
free basis. When both $M$ and $P$ occur, we specify in which group the
translation is primitive.

\subsection{Uniform geometric estimates}

Write $\mathcal C_Q$ for the unit-edge Cayley tree of a finitely
generated free group $Q$ with a fixed basis. A homomorphism on vertices
is extended over each edge by its reduced image path. The image of a
geodesic under a fixed injective homomorphism is a uniform
quasigeodesic; its tightened line is at uniformly bounded Hausdorff
distance from it. We first record the consequence needed both here and
in the final argument.

\begin{lemma}[Nearness and order under a fixed transfer]
\label{exp:transport}
Let $a:Q\to F$ and $b:Q\to F'$ be injective homomorphisms between
finitely generated free groups with fixed bases. For every $D\ge0$
there are constants $E,T\ge0$ with the following properties.
Let $L$ be an oriented line of $\mathcal C_Q$, and let $L_a,L_b$
be its tightened image lines in $\mathcal C_F,\mathcal C_{F'}$.
If $q\in Q$ satisfies $d(a(q),L_a)\le D$, then
\[
 d(b(q),L_b)\le E.
\]
If two such elements $q_1,q_2$ have projections to $L_a$ positively
ordered and separated by more than $T$, then their projections to
$L_b$ are in the same order. As their separation on $L_a$ tends to
infinity, their separation on $L_b$ tends to infinity, uniformly over
all $L,q_1,q_2$.
\end{lemma}

\begin{proof}
Fix Hausdorff bounds between each tightened image line and the images
of vertices of $L$. Choose a vertex $u\in L$ whose $a$-image is
uniformly close to a point of $L_a$ closest to $a(q)$. The lower
quasi-isometry bound for $a$ gives a bound on $d_Q(q,u)$ depending
only on $D$ and $a$. The upper bound for $b$, followed by its
line Hausdorff bound, proves the first assertion.

For the order assertion, first consider actual vertices $u<v$ on an
oriented source line and one fixed embedding $f$. The image of $v$ is
within a fixed distance of the geodesic from $f(u)$ to $f(z)$ for every
source vertex $z$ sufficiently far beyond $v$. Projecting to the
tightened image line and letting $z$ tend to its positive end shows
that the projection of $f(v)$ cannot lie more than a fixed distance
behind that of $f(u)$. On the other hand, the distance between these
projections is at least
\[
 c_f d_Q(u,v)-C_f
\]
for fixed $c_f>0,C_f\ge0$, by the lower embedding bound and the
line Hausdorff bound. Thus sufficiently separated vertices have
projections in the correct order, with separation tending to infinity.

Apply this first to $a$ and then to $b$. Replace $q_i$ by vertices
$u_i\in L$ at the uniformly bounded distance obtained in the first
paragraph. Their image projections change by uniformly bounded
amounts. Sufficient positive separation of the $a(q_i)$ therefore
forces $u_1<u_2$ with large separation, and then forces the asserted
order and separation of the $b(q_i)$. All constants came from the
fixed embeddings and $D$.
\end{proof}

We will also use two elementary properties of trees. If $x,y$ lie
within $D$ of a geodesic line $L$, the part of $[x,y]$ farther than
$D$ from both endpoints lies on $L$. Indeed, project $x,y$ to $L$;
if the projections are distinct, the connecting segment enters and
leaves $L$ there. If they coincide, then $d(x,y)\le2D$, so there
is no point farther than $D$ from both endpoints.
The same assertion holds for an arbitrary convex subtree. Secondly,
the lines meeting a fixed finite ball form a compact subset of line
space: they are a finite union of the compact sets of lines through
its vertices. A current has finite mass on this set.

\begin{lemma}[Expansion with fixed trimming]
\label{exp:trimming}
Fix a maximal carrier $M$ and the maps above. There are constants
$Q\ge0$ and $L_*>0$ with the following property. Let $L$ be an
oriented $M$-Cayley-tree line, let $x,y$ be vertices on $L$ with
$d_M(x,y)=l\ge L_*$, and let $\mathcal A\subset\mathcal U$ be
a Borel set whose lines lie within $Q$ of both $x,y$ in
$\mathcal C_P$. Suppose a sign $\epsilon\in\{+1,-1\}$ satisfies
\[
 d_M(\beta^\epsilon x,\beta^\epsilon y)\ge\lambda l.
\]
If every line of $\mathcal A$ belongs to
$\partial^2\Lambda S_\epsilon$, then its image
$\mathcal A'=\Theta_\epsilon(\mathcal A)$ has the same measure,
lies in $\mathcal U$, and admits vertices $x',y'$ on the tightened
line $\beta^\epsilon L$ such that
\[
 d_P(x',\ell),d_P(y',\ell)\le Q\quad(\ell\in\mathcal A'),
 \qquad d_M(x',y')\ge2l.
\]
The same sign satisfies
\[
 d_M(\beta^\epsilon x',\beta^\epsilon y')
 \ge\lambda d_M(x',y').
\]
The constants are independent of $L,x,y,\mathcal A$ and of the
number of times this assertion is applied.
\end{lemma}

\begin{proof}
Choose $Q$ bounding the distance of every $M$-vertex on a geodesic
segment or line to its tightened $P$-image, and bounding the reverse
distance from that image to images of vertices. Increase it if
necessary. Applying Lemma~\ref{exp:transport} to the two fixed
window maps gives a common $C_1$ such that
$u=\beta^\epsilon x$ and $z=\beta^\epsilon y$ are within $C_1$
of every line in $\mathcal A'$.

Choose lower embedding constants $a>0,b\ge0$ with
\[
 d_P(g,h)\ge a d_M(g,h)-b\qquad(g,h\in M).
\]
Let $B$ be a common vertex-to-line bound for the two automorphisms
$\beta^{\pm1}$ of $M$. Choose an integer $R$ such that
\begin{equation}
 R>B,\qquad aR-b-Q>C_1.
 \label{exp:R-choice}
\end{equation}
Delete $R$ edges at both ends of $[u,z]$ in $\mathcal C_M$, and
call its new endpoints $x',y'$. The first inequality in
\eqref{exp:R-choice} puts them on the tightened line
$\beta^\epsilon L$. Their $P$-images lie within $Q$ of
$[u,z]_P$. If $v$ is either new endpoint, its projection to this
segment has distance at least $aR-b-Q>C_1$ from each end.
The intervening part of $[u,z]_P$ lies on every line of
$\mathcal A'$, by the tree observation. Hence $x',y'$ are still
within $Q$ of every such line.

Put
\[
 K=\max\{\operatorname{Lip}(\beta),
                 \operatorname{Lip}(\beta^{-1})\}.
\]
The new length $l'$ satisfies
\begin{equation}
 l'\ge\lambda l-2R,
 \qquad
 d_M(\beta^{-\epsilon}x',\beta^{-\epsilon}y')
 \le l+2KR.
 \label{exp:trim-estimates}
\end{equation}
Choose $L_*$ sufficiently large that, for every $l\ge L_*$,
\[
 \lambda l-2R\ge2l,
 \qquad l+2KR<2\lambda l.
\]
Then the opposite direction at the new pair has length strictly
less than $\lambda l'$. Inequality~\eqref{exp:based-inequality}
forces expansion in the original direction. Measure preservation
and membership in $\mathcal U$ follow from the invariant conull
transfer construction. Since the maps are fixed boundary
homeomorphisms on their domains, the image set is Borel.
\end{proof}

\subsection{A failed return produces a periodic segment}

Fix the carrier $M$, its return maps, and the constants in
Lemma~\ref{exp:trimming}. The next lemma supplies the common construction
for both support arguments. Its reference line is transferred in the
boundary of $M$; only the competing line in its conclusion is asserted
to lie in $\Omega$.

\begin{lemma}[Finite failure of an expanding return]
\label{exp:finite-failure}
There are constants $L_0\ge L_*$, $c>0$, $C_2\ge0$, and an integer
$D\ge1$, depending only on the fixed splitting, carrier, and return maps, with
the following property. Let $\eta\in\Omega\cap\partial^2\Lambda M$,
let $L$ be its corresponding oriented $M$-Cayley-tree line, and choose
vertices $x_0<y_0$ on $L$ with $l_0=d_M(x_0,y_0)\ge L_0$. Suppose
$\epsilon\in\{+1,-1\}$ satisfies
\[
 d_M(\beta^\epsilon x_0,\beta^\epsilon y_0)\ge\lambda l_0.
\]
Then there are $n\ge0$, a line $\ell\in\mathcal U$, and an oriented
periodic $P$-line $\zeta$ of primitive period at most $D$ such that
$\ell$, $\zeta$, and the tightened $P$-image $\eta_n$ of
$\beta^{\epsilon n}L$ share a segment of length at least
\begin{equation}
 c\,2^n l_0-C_2.
 \label{exp:failure-overlap}
\end{equation}
The endpoints of $\zeta$ lie in the limit sets of the stabilizers of
two distinct actual windows of the shape defining $\Theta_\epsilon$.
One is the prescribed window, with stabilizer $S_\epsilon$.
\end{lemma}

\begin{proof}
Choose $L_0\ge L_*$ large enough that the projections of $x_0,y_0$
to $\eta$ are distinct and occur in that order. This choice is uniform
by Lemma~\ref{exp:transport}; we will increase $L_0$ once more below.
The cylinder prescribing the directed segment between those projections
has positive finite measure, since $\eta\in\Omega$. Its intersection
with $\mathcal U$ is a Borel set $\mathcal A_0$ of positive finite
measure, every line of which is within $Q$ of $x_0,y_0$.

Attempt to apply $\Theta_\epsilon$ successively to the entire
current set. At time $n$, call the attempt a \emph{failure} if
$\mathcal A_n$ is not contained in
$\partial^2\Lambda S_\epsilon$. Otherwise define
$\mathcal A_{n+1}=\Theta_\epsilon(\mathcal A_n)$ and use
Lemma~\ref{exp:trimming} to obtain the next segment. As long as
the process continues, its vertices $x_n,y_n$ lie on
$\beta^{\epsilon n}L$, their distance $l_n$ is at least
$2^n l_0$, and every line of $\mathcal A_n$ lies within $Q$ of
both vertices. Its measure is the fixed number
$\delta=\mu(\mathcal A_0)>0$.

\smallskip\noindent
\emph{The process cannot continue forever.}
Choose an $M$-vertex $a_n$ within one edge of the midpoint of
$[x_n,y_n]_M$ and
translate $\mathcal A_n$ by $a_n^{-1}$. The central portion of
$[x_n,y_n]_P$ belongs to every line of $\mathcal A_n$.
The image of $a_n$ is uniformly close to this segment, and the
lower embedding bound puts its projection arbitrarily far from
both endpoints as $n$ increases. That projection therefore lies
in the common central portion. Thus the translated sets
eventually lie in one compact set $\mathcal K$ of lines meeting a
fixed ball. Their measures are still $\delta$.

Every subsequential limit of these lines belongs to
$\partial^2\Lambda M$. To see this, take a subsequence of the
centered $M$-segments. Both halves tend to infinite length, so a
further subsequence converges to a bi-infinite $M$-geodesic. Its
quasigeodesic image in $P$ has two distinct boundary endpoints,
and its tightened line is the limit of the growing common central
segments. These endpoints lie in $\Lambda M$.

Consequently every relative neighborhood of
$\mathcal K\cap\partial^2\Lambda M$ in $\mathcal K$ contains
all sufficiently late translated sets. Otherwise points outside
one such neighborhood would have a contrary subsequential limit
in the compact set $\mathcal K$. Outer regularity of the finite
measure on $\mathcal K$ now gives
\[
 \mu(\mathcal K\cap\partial^2\Lambda M)\ge\delta,
\]
contrary to \eqref{exp:zero-carrier-mass}.

\smallskip\noindent
\emph{A failure produces a bounded-period axis.}
Suppose failure occurs at time $n$, and choose
$\ell\in\mathcal A_n\setminus\partial^2\Lambda S_\epsilon$.
This line has an actual trajectory for the same finite one-sided
window. Its stabilizer $K$ is noncyclic, and this actual path differs
from the prescribed path: equality would put both endpoints of
$\ell$ in $\Lambda S_\epsilon$.

Both $\ell$ and the tightened image $\eta_n$ of
$\beta^{\epsilon n}L$ pass within $Q$ of $x_n,y_n$. Since
$x_n,y_n\in M\le S_\epsilon$, their distances from
$\operatorname{Hull}(\Lambda S_\epsilon)$ equal the fixed number
$d_P(1,\operatorname{Hull}(\Lambda S_\epsilon))$. Set
\[
 H_\epsilon=\max\{Q,
             d_P(1,\operatorname{Hull}(\Lambda S_\epsilon))\}.
\]
Deleting $H_\epsilon$ from both ends of $[x_n,y_n]_P$ leaves a
segment in $\ell$, in $\eta_n$, and in the prescribed stabilizer
hull. It also lies in $\operatorname{Hull}(\Lambda K)$ because
$\ell$ does. Its length is at least
\begin{equation}
 a l_n-C,
 \label{exp:central-overlap}
\end{equation}
where $a>0,C\ge0$ depend only on the fixed maps and embeddings.
In particular, no constant from the varying competing subgroup $K$
is needed.

Read this segment in the two finite stabilizer cores. Their fiber
product pairs core states, with an edge for each pair of equally
labeled edges. At an actual vertex $z$ the corresponding pair of states represents the two
paths translated by $z^{-1}$. They remain distinct. Every loop in
the corresponding component of the fiber product therefore fixes
both paths, so its fundamental group injects into an at-most-cyclic
group, by Proposition~\ref{stat:windows}. The component has rank at most
one. Since the window is fixed, there are only finitely many core
representatives, and the numbers of unoriented edges in all relevant fiber products
are bounded by one integer $D\ge1$. Conjugating a stabilizer preserves its labeled
unbased core and changes only the state used for this reading.

Increase $L_0$ so that $aL_0-C>2D+1$, taking the maximum of
the constants for the two signs. This ensures that every failure
reading is long enough to contain a circuit. A reduced path in a
finite tree has length at most its number of edges. A connected
rank-one finite graph consists of one circuit with trees attached,
so a sufficiently long reduced path follows that circuit between
initial and terminal pieces of length at most $D$.
Our reading therefore has a central periodic portion of length at
least $a l_n-C-2D$. Extend it in the absolute Cayley tree to an
oriented periodic line $\zeta$. Its primitive $P$-period is at most
$D$. Unfolding the circuit in each core puts its endpoints in both
actual window stabilizer limit sets. Since $l_n\ge2^n l_0$, this
proves \eqref{exp:failure-overlap}, with constants enlarged uniformly
over the two signs. The line $\ell$ belongs to $\mathcal U$, as
required.
\end{proof}

The failure need not occur on a set of positive measure. Its chosen
line already belongs to $\mathcal U\subset\Omega$, which is enough
for the support limit below. Positive mass rules out indefinite
transfer of the whole cylinder. Figure~\ref{exp:failure-figure}
shows how the two actual windows produce the periodic middle.

\begin{figure}[ht]
\centering
\begin{tikzpicture}[every node/.style={font=\small}]
\draw[gray,thick] (0,2.4) -- (10,2.4);
\draw[gray,thick] (0,3.2) -- (2,2.4);
\draw[gray,thick] (8,2.4) -- (10,3.2);
\draw[blue!55!black,very thick] (2,2.4) -- (8,2.4);
\draw[blue!55!black,densely dashed,thick] (.8,1.6) -- (2,2.4);
\draw[blue!55!black,densely dashed,thick] (8,2.4) -- (9.2,1.6);
\node[left] at (0,2.4) {$\eta_n$};
\node[left] at (0,3.2) {$\ell\in\Omega$};
\node[left,blue!55!black] at (.8,1.6) {$\zeta$};
\node[above] at (5,2.55) {periodic middle of the common segment};
\draw[gray,thick] (1.7,0) -- (4.4,0);
\draw[gray,thick] (5.6,0) -- (8.3,0);
\draw[blue!55!black,very thick] (5,0) circle (.6);
\fill (4.4,0) circle (1.7pt);
\fill (5.6,0) circle (1.7pt);
\draw[-{Stealth[length=2mm]},thin] (5,.75) -- (5,2.05);
\node[right,align=left] at (5.15,1.35) {unfold the\\circuit};
\node[below] at (3,-.1) {at most $D$ edges};
\node[below] at (7,-.1) {at most $D$ edges};
\node[below] at (5,-.65) {circuit length at most $D$};
\node[below] at (5,-1.3)
 {a rank-one component of $C(S_\epsilon)\mathbin{\times_R}C(K)$};
\end{tikzpicture}
\caption{At a failed return, $\eta_n$ and a competing support line
$\ell$ share a long segment in two window-stabilizer hulls. Its reading
in the core fiber product lies in a rank-one component, shown below.
The initial and terminal tree portions have bounded length; between
them the reading follows the circuit. Unfolding that circuit gives the
periodic middle and its extension $\zeta$. The figure is schematic;
only $\ell$ is asserted to lie in the support.}
\label{exp:failure-figure}
\end{figure}

\begin{proposition}
\label{exp:carrier-periodic}
If $\Omega$ contains a line carried, up to translation, by a maximal
noncyclic carrier $M$, then $\Omega$ contains a periodic line.
\end{proposition}

\begin{proof}
Translate the given line into $\partial^2\Lambda M$. Choose
successively longer segments on its corresponding $M$-line and,
for each segment, a sign supplied by \eqref{exp:based-inequality}.
Apply Lemma~\ref{exp:finite-failure}, passing to a subsequence with
the same sign. Its support line and bounded-period axis share a
segment whose length tends to infinity.

Recenter this shared segment
by an $M$-translation chosen near its midpoint on $\eta_n$.
The translated periodic axes have period at most $D$ and meet a
fixed ball. There are only finitely many such axes: there are only
finitely many cyclic labels of length at most $D$, and finitely many
ways for one of their axes to meet that ball. Along a subsequence
the periodic axis is fixed. The translated lines $\ell\in\Omega$
agree with it on segments tending to infinity in both directions,
so they converge to it. Closedness of $\Omega$ proves that this
periodic line belongs to $\Omega$.
\end{proof}

\subsection{Pulling back whole periods}

The next lemma controls the number of periods in a common segment.
Controlling only its metric length would be insufficient, since a
period can grow under an inverse automorphism.

\begin{lemma}[Bounded loss of periods]
\label{exp:period-pullback}
Let $\psi$ be an automorphism of a finitely generated free group $M$
with a fixed basis. There is an integer $C_4\ge0$, depending only on
$\psi$, with the following property. Suppose two lines $L_1,L_2$
of $\mathcal C_M$ share a segment $[u,z^k u]$, where $u$ is a
vertex, $z$ translates $L_2$, and $k>C_4$. Their tightened
$\psi$-images share a segment consisting of $k-C_4$ whole
$\psi(z)$-periods. If $z$ is primitive as a translation in $M$,
then so is $\psi(z)$.
\end{lemma}

\begin{proof}
Choose an integer $B\ge0$ such that the image of any vertex on any
line is within $B$ of the tightened image line. Such a uniform
constant exists by the quasigeodesic line bound for the fixed map
$\psi$. Put $q=2B+1$ and $C_4=2q$.

Write $L_i'=\psi L_i$ for the tightened images and $z'=\psi(z)$.
The line $L_2'$ is the axis of $z'$. If $a$ is the projection of
$\psi(u)$ onto this axis, the projection of the other endpoint
image is exactly
\[
 \pi_{L_2'}\bigl(\psi(z^ku)\bigr)
 =\pi_{L_2'}\bigl(z'^k\psi(u)\bigr)=z'^k a,
\]
because projection onto an invariant line commutes with its
translation. Both endpoint images lie within $B$ of both output
lines. Thus $a,z'^ka$ lie within $2B$ of $L_1'$.

The translation length of $z'$ is an integer at least one. The
vertices $z'^q a$ and $z'^{k-q}a$ therefore lie farther than
$2B$ from both ends of $[a,z'^ka]$ and occur in the stated order.
The tree observation puts their joining segment on $L_1'$, and
it already lies on $L_2'$. It has exactly $k-2q=k-C_4$ full
periods, with both endpoints in the same integer $z'$-orbit.
There is no additional phase loss. Automorphisms preserve being a
non-power, proving the final assertion.
\end{proof}

\begin{lemma}[Agreement of periodic labels]
\label{exp:periodic-labels}
Two bi-infinite label sequences with respective positive integer
periods $a,b$ that agree on $a+b$ consecutive positions agree
everywhere. Consequently two oriented periodic lines in a labeled
Cayley tree with such a common segment are the same actual line.
\end{lemma}

\begin{proof}
Number the common positions $0,\ldots,a+b-1$. For
$0\le i<a$, the letter at position $i$ equals the letter at
$i+b$ by the second period. By the first period this equals the
letter at $(i+b)\bmod a$. Thus the first period pattern is
invariant under addition by $b$ modulo $a$, and is periodic with
period $d=\gcd(a,b)$. The first $b$ common letters then show that
the second pattern is the same $d$-periodic pattern. Both entire
sequences therefore agree. A label and an initial vertex determine
a unique path in a Cayley tree, so equality of the sequences and
the shared segment identifies the actual lines.
\end{proof}

\begin{lemma}[A bounded-period axis near a carrier]
\label{exp:carrier-period}
Let $M\le P$ be a finitely generated nontrivial subgroup, with fixed
bases, and fix $D\ge1$. There are constants $L_D,B_D$ such that if
a periodic $P$-line $\zeta$ of primitive period at most $D$ shares
a segment of length at least $L_D$ with an $M$-line, then $\zeta$
is itself an $M$-line. Its primitive translation in $M$ has cyclic
$M$-length at most $B_D$.
\end{lemma}

\begin{proof}
Let $V$ be the number of vertices in the finite core of $M$ in the
$P$-rose, and take $L_D=(V+1)D$. The common segment lies in
$\operatorname{Hull}(\Lambda M)$. If $r$ is the actual primitive
$P$-translation of $\zeta$, that segment contains $V+1$ successive
vertices $r^i z_0$ in one period phase. Two have the same core state:
\[
 M r^i z_0=M r^j z_0\qquad(0\le i<j\le V).
\]
Cancelling $z_0$ gives $r^{j-i}\in M$. Thus a power of the
translation of this very axis belongs to $M$, and the whole line
has endpoints in $\Lambda M$.

If $z$ is its primitive translation in $M$, it is the smallest
positive power of $r$ belonging to $M$, up to orientation. Hence
$z=r^q$ for some $1\le q\le V$. Choose embedding constants with
\[
 |g|_M\le A|g|_P+B\qquad(g\in M).
\]
Apply this to powers of $z$, divide by the exponent, and take the
limit to obtain
\begin{equation}
 \ell_M(z)\le A\ell_P(z)\le AVD.
 \label{exp:carrier-period-bound}
\end{equation}
Thus $B_D=AVD$ suffices. The two primitive translations need not
coincide, which is why both groups are specified.
\end{proof}

\subsection{Exclusion of periodic support lines}

\begin{proposition}
\label{exp:no-periodic}
The support $\Omega$ contains no periodic line.
\end{proposition}

\begin{proof}
Suppose $\eta_0\in\Omega$ is periodic.

\smallskip\noindent
\emph{Step 1: construct a trajectory and locate a maximal carrier.}
We first construct one compatible trajectory for $\eta_0$; uniqueness
at its rational endpoints is not assumed.

Approximate $\eta_0$ by lines of $\mathcal U$. For each fixed
centered window, their stabilizer hulls meet a common finite ball.
Lemma~\ref{stat:window-finiteness} lets us pass to a subsequence
with a constant actual window path. Taking these subsequences
successively gives a bilateral aligned path whose every finite
window has noncyclic stabilizer. The limiting endpoints belong to
each such stabilizer limit set. Choose frames $h_s$ on this path,
with $h_0=1$. Boundary transfer on a fixed window is continuous, so
the corresponding conull image lines converge in the successive
frames to lines $\eta_s\in\Omega$. They are periodic: the
core-state repetition argument in Lemma~\ref{exp:carrier-period}
puts a power of the original
axis translation in each fixed window stabilizer, and this power
is transferred to a nontrivial element translating the image line.

Choose the primitive $P$-translation $r_s$ of $\eta_s$ toward its
positive endpoint. Write a forward step as
$h_{s+1}=h_s p_s t q_s$, with $p_s,q_s\in P$.
There is a least positive integer $m_s$ with
$r_s^{m_s}\in p_sIp_s^{-1}$. Its image is
$r_{s+1}^{n_s}$, where $n_s$ is the least positive integer with
$r_{s+1}^{n_s}\in q_s^{-1}Jq_s$. Minimality on the second side
follows by applying the inverse subgroup isomorphism to its full
axis stabilizer. The exponents are positive because we oriented
each translation toward the transferred positive endpoint.

Recall the finite ambiguity dictionary of
Lemma~\ref{stat:finite-roots}: if either minimum exponent exceeds
one, or if one of
these axes belongs to two distinct domain cosets on its side,
one of its primitive root classes lies in a fixed finite set, up
to conjugacy and inversion. Infinitely many such indices would
repeat a root class along this actual reduced path.
Lemma~\ref{stat:repetition} would give a tree-hyperbolic element
commensurating a nontrivial elliptic cyclic group, contrary to
BS-freeness. Therefore only finitely many indices in the entire
bilateral trajectory are exceptional in either way.

It follows that one can choose positive integers $k_s$ for all
$s\in\mathbb Z$ with
\begin{equation}
 m_s\mid k_s,\qquad k_{s+1}=n_s k_s/m_s.
 \label{exp:resident-exponents}
\end{equation}
Only finitely many ratios differ from one. Choose $k_0$ divisible
by all the finitely many factors needed when propagating these
relations in both directions. Then
\[
 h_s r_s^{k_s}h_s^{-1}
\]
is one fixed nontrivial element, independent of $s$.

Fix $s$ temporarily. The cyclically tightened circuit representing
$r_s^{k_s}$ lifts closed to
a component of each centered window core $\Gamma_n$, after
recentering the trajectory at
$h_se$: the edges $E_j'=h_s^{-1}E_{s+j}$ have $E_0'=e$ and so
belong to the same family of windows defining $\Gamma_n$. Its length is fixed as the window size tends to infinity.
It cannot traverse a natural edge whose length tends to infinity:
any immersed closed path using the interior of such an edge
traverses the whole edge, cyclically. Thus $\eta_s$ is carried by
the fixed small subgraph. If infinitely many positive times were
carried only by cyclic small components, their primitive
$P$-roots would again belong to finitely many conjugacy classes
and repeat. We may therefore choose a time beyond every
exceptional forward index at which the periodic line is carried
by a noncyclic small subgroup. Enlarge it to a maximal carrier
$M$ and translate the line into $\partial^2\Lambda M$.
If there were no maximal noncyclic carrier, this conclusion itself
would be a contradiction.

\smallskip\noindent
\emph{Step 2: follow expanding return iterates.}
The subsequent forward path agrees with the path fixed by $M$.
At its first edge, both choices allow the periodic endpoints,
and the domain coset is unique at this late time. Transfer and
repeat to get agreement at every forward edge. In particular all
positive $\beta$-iterates of the line belong to $\Omega$ and have
unique forward choices along the prescribed return windows.
This conclusion comes from the constructed trajectory, not from
an as yet unproved support-invariance assertion.

Let $a_j\ge1$ be the primitive cyclic $M$-length of its $j$th
positive $\beta$-iterate. At some time one has
$a_{j+1}\ge\lambda a_j$. Otherwise the cyclic version of
\eqref{exp:based-inequality}, applied at every $j\ge1$, would
give $a_{j-1}\ge\lambda a_j$, impossible for positive integer
lengths. Once a forward expansion occurs, it continues: the
backward neighbor is too short to supply the next expanding
alternative. Restart at such a time, and call the initial
primitive period length $d_0$.

Let $z$ be the positive primitive $M$-translation of this initial line.
Choose a vertex $x_0$ on its axis and $y_0=z^k x_0$, where $k$ is arbitrarily
large. Then $l_0=kd_0$. The cyclic expansion gives
$\ell_M(\beta z)\ge\lambda d_0$, and therefore
\[
 d_M(\beta x_0,\beta y_0)
 =d_M(\beta x_0,(\beta z)^k\beta x_0)
 \ge k\ell_M(\beta z)\ge\lambda l_0.
\]
This verifies the based-segment hypothesis even if $\beta x_0$
is off the axis of $\beta z$.
For $l_0\ge L_0$, Lemma~\ref{exp:finite-failure} now supplies a
failure time $n$, a support line, and a bounded-period axis sharing
\eqref{exp:failure-overlap} with the $n$th return image.

\smallskip\noindent
\emph{Step 3: pull a failure axis back to the initial period.}
The resulting periodic $P$-line $\zeta$ has uniformly bounded period,
allows both the prescribed and a competing path, and shares with the current $M$-line a segment of length at least
\eqref{exp:failure-overlap}.
Take $l_0$ large enough that \eqref{exp:failure-overlap} exceeds
the fixed threshold in Lemma~\ref{exp:carrier-period}, already at
$n=0$. That lemma implies
that $\zeta$ is itself an $M$-line, with primitive
$M$-period bounded by \eqref{exp:carrier-period-bound}.

The two corresponding $M$-lines share a segment of length at
least $c_0 2^n l_0-C_0$, for fixed $c_0>0,C_0\ge0$.
Here are the details of this change of metric. Near each end of
their common $P$-segment, choose a vertex on each $M$-line whose
$P$-image is within the fixed line Hausdorff bound. At each end
the two chosen vertices are uniformly close in $M$, by the
lower embedding estimate. The vertices chosen near opposite
ends are separated in $M$ by a positive multiple of the
$P$-length minus a constant, by the upper embedding estimate.
Choose the pairs in their oriented order, as justified by
Lemma~\ref{exp:transport}. In a tree, the two joining segments
therefore have a common central segment after deleting a fixed
amount at both ends. This gives the asserted lower bound.

The uniformly bounded primitive period of $\zeta$ now shows
that this common segment contains at least
\begin{equation}
 c'2^n l_0-C_3
 \label{exp:initial-period-count}
\end{equation}
whole primitive $\zeta$-periods, for fixed $c'>0,C_3\ge0$.
Apply Lemma~\ref{exp:period-pullback} with $\psi=\beta^{-1}$
successively $n$ times. At time zero at least
\begin{equation}
 c'2^n l_0-C_3-nC_4
 \label{exp:accumulated-period-loss}
\end{equation}
whole periods remain. Choose the initial length so large that
\begin{equation}
 c'l_0>\max\{C_4,\ C_3+d_0+1\}.
 \label{exp:uniform-in-time}
\end{equation}
Since $2^n\ge n+1$, expression
\eqref{exp:accumulated-period-loss} exceeds $d_0+1$ for every
possible failure time $n$. In particular all intermediate
applications of the pullback lemma have enough periods.

The primitive period of the second line after these pullbacks
may be large; call it $b\ge1$. The number of full periods,
rather than its metric length, was controlled. More than
$d_0+1$ periods give a common label segment of length at least
$(d_0+2)b\ge d_0+b$. Lemma~\ref{exp:periodic-labels} identifies
the original periodic line and the pulled-back $\zeta$ as the
same actual $M$-line. Applying $\beta^n$ identifies their
lines at failure, in both the $M$- and $P$-trees. This line
allows both the prescribed and the competing forward window.
That contradicts the unique forward coset choices on the late
periodic trajectory. Failure is impossible as well, completing
the contradiction.
\end{proof}

\subsection{The trajectory construction on the entire support}

We have excluded periodic support points without assuming that all
support points admit unique transfers. We can now prove precisely
that extension, together with the uniform bound needed to count
special rays.

\begin{corollary}
\label{exp:support}
The following assertions hold for the closed support $\Omega$.
\begin{enumerate}
 \item No line of $\Omega$ is carried by the small subgraph.
 There is an integer $B_0$ such that every segment of a lift of
 an $\Omega$-line to any of the selected subsequence cores
 $\Gamma_n$, if contained entirely in the small subgraph, has
 length at most $B_0$.
 \item Every endpoint of every line of $\Omega$ is irrational.
 \item Every line of $\Omega$ lifts to each centered core union $\Gamma_n$
 and has a unique compatible bilateral trajectory. On each
 fixed window, the partial boundary transfer takes every
 support line in its domain to a line of $\Omega$.
\end{enumerate}
\end{corollary}

\begin{proof}
An immersed bi-infinite path in a finite cyclic small component
is periodic, and such a line is excluded by
Proposition~\ref{exp:no-periodic}. A line carried by a noncyclic
small component is carried, up to conjugacy, by a maximal
noncyclic carrier. Proposition~\ref{exp:carrier-periodic} would
then produce a periodic support line, also impossible. A tree
component cannot carry an immersed bi-infinite path. This proves
the first assertion about whole lines.

If small-only segments of lifted support lines had unbounded
length, center progressively longer such segments at an actual
vertex and translate that vertex to the identity in $P$.
The translated lines remain in $\Omega$ and lie in the compact
set of lines through the identity. Since the small graph is
fixed and finite, a subsequence of the centered lifted segments
converges to an immersed bi-infinite path in it. The corresponding
limit support line would be carried by the small graph. This
contradiction gives the uniform bound $B_0$.

A rational endpoint is represented by an eventually periodic
ray. If such a ray lay on a support line, translating farther
and farther along its periodic tail would give a periodic line
as a limit of support lines. Closedness of $\Omega$ and
Proposition~\ref{exp:no-periodic} exclude this. Both endpoints
of every support line are therefore irrational.

The existence of lifts to each finite core was already proved
in Lemma~\ref{car:window-cores}: it follows from conull lifts,
density, and closedness of liftability into a fixed finite
graph. Each such lift gives an actual window path carrying the
line's endpoints. The irrational-endpoint uniqueness in
Proposition~\ref{stat:windows} makes these paths compatible when a
window is shortened. They therefore give a unique bilateral
trajectory.

It remains to check the support of transferred lines. Approximate
a given support line by lines of $\mathcal U$ and fix one
window. All approximants meet a common fixed ball, so the hulls
of their window stabilizers do too. There are only finitely many
possible actual paths. Any path occurring infinitely often
carries the limiting endpoints, because its limit set is
closed; irrational uniqueness identifies it with the limiting
line's own path. Thus the approximants eventually have that
same path. The boundary transfer on its fixed subgroup is
continuous. Their images belong to $\mathcal U\subset\Omega$
and converge to the transferred limiting line, which is in
$\Omega$ by closedness. This proves the final assertion without
discarding any support points.
\end{proof}
